
    \documentclass[a4paper]{article}
    \usepackage{savetrees}
    \usepackage{hyperref}
    \usepackage{pdfpages}

    \begin{document}

    Ameya Godbole and Robin Jia.  {\bf \href{http://cs.umd.edu/~jbg//docs/2025_aclfindings_verify.pdf}{Verify with Caution: The Pitfalls of Relying on Imperfect Factuality Metrics}}.  \emph{Findings of the Association for Computational Linguistics}, 2025.



    \vspace{5cm}

    \begin{verbatim}
    @article{Godbole:Jia-2025,
	Title = {Verify with Caution: The Pitfalls of Relying on Imperfect Factuality Metrics},
	Author = {Ameya Godbole and Robin Jia},
	Journal = {Findings of the Association for Computational Linguistics},
	Year = {2025},
	Url = {http://cs.umd.edu/~jbg//docs/2025_aclfindings_verify.pdf},
}

    \end{verbatim}

    {\bf Accessible Abstract:} To rigorously evaluate language model outputs, we must identify reliable evaluation metrics for long-form question answering, which is a longstanding challenge in natural language processing. Recently, there has been growing optimism that LLMs themselves can serve as reliable evaluators. We systematically studied the reliability of these LLM-based metrics across 11 natural language generation datasets. We uncovered numerous weaknesses when these methods are deployed naively: they are mutually inconsistent, misestimate system-level performance, and are overly reliant on simple heuristics. Moreover, we propose a metric adjustment technique that reduces system-level bias of these metrics, enabling more reliable evaluation. Overall, we establish best practices for LLM-based evaluation that will enable us to conduct more reliable evaluation of our future question answering systems.\\~

    

    

    \vspace{3cm}
    Downloaded from \url{http://cs.umd.edu/~jbg/docs/2025_aclfindings_verify.pdf}  \\

    \vspace{2cm}
    \textit{Contact Jordan Boyd-Graber (jbg@boydgraber.org) for questions about this paper.}

    \includepdf[pages={-}]{src_docs/2025_aclfindings_verify.pdf}

    \end{document}
    